\documentclass[../main]{subfiles}
\begin{document}
\section{Acuerdos áulicos}
\subsection*{Normas para la presentación de la carpeta de trabajos prácticos, teoría y acuerdos áulicos.}
\begin{enumerate}
    \item La portada debe ser la propuesta por el docente, en caso que se encuentra la portada de la escuela se utilizará la misma, el índice debe estar completo, con la numeración completa, en la misma debe estar el programa de la materia y la presentes pautas que deben estar firmadas por el alumno.
    \item Rótulos completos, realizados en regla, con todos los datos correspondientes y el formato propuesto por el docente, las hojas deben estar numeradas, letra técnica o imprenta prolija segun lo pactado por el docente.
    \item La carpeta debe estar prolija, se tendrá en cuenta la ortografía, las hojas deben ser la misma sen cada trabajo práctico, al cambiar de trabajo pueden cambiar las hojas, se debe utilizar el mismo color de lapicera, los dibujos deben ser realizados en regla.
    \item El alumno deberá tener dos evaluaciones por informe, en total seis evaluaciones al año, si el alumno falta deberá traer el certificado correspondiente.
    \item \textbf{Trabajo en clase, trabajo en grupo, pizarrón, actitudinal:} El alumno deberá tener un vocabulario adecuado en el aula, respeto hacia el docente y hacia sus compañeros y las normas de convivencia.
    \item Carpeta completa y entrega en termino.
    \item Asistencia (80\%).
\end{enumerate}
\vspace{1cm}
\begin{center}
Firma Alumno $\dots\dots\dots\dots\dots\dots\dots\dots\dots\dots$
\end{center}

\end{document}
